% Battery Inspection System report sections for Overleaf
% Upload this file to the Overleaf project and insert it with:
% % Battery Inspection System report sections for Overleaf
% Upload this file to the Overleaf project and insert it with:
% % Battery Inspection System report sections for Overleaf
% Upload this file to the Overleaf project and insert it with:
% % Battery Inspection System report sections for Overleaf
% Upload this file to the Overleaf project and insert it with:
% \input{Overleaf_Sections_3_2_4_to_3_2_7}
%
% These headings use \subsubsection so that, when placed under Section 3.2,
% LaTeX can number them automatically as 3.2.4--3.2.7.
% Add \label{fig:dofbot-robot-camera-control} to the Figure 7 environment.

\subsubsection{System Layout}
\label{sec:system-layout}

The proposed battery inspection and sorting system is organized into three main processing areas: the Windows inspection station, the conveyor section, and the Jetson--DOFBOT robot station. This distributed arrangement separates image classification, object transportation, trigger detection, and robot motion control into independent modules. The modular design simplifies testing and allows each subsystem to be developed or replaced without redesigning the entire system.

The Windows inspection station consists of a Windows computer or mini workstation, an IMITECH IMB-770GC industrial camera, the MVS camera driver, the WPF human--machine interface, and the YOLOv8 inference component. The IMITECH camera is positioned above the conveyor to capture detailed images of each battery. Its primary function is to provide high-quality input images for the AI model. Camera communication is handled through the MVS SDK rather than a conventional webcam interface. Therefore, the camera stream must not be opened simultaneously by both the MVS application and the WPF software.

The YOLOv8 model classifies each inspected battery into one of four categories:

\begin{itemize}
    \item \texttt{dented}: the battery casing is deformed or dented;
    \item \texttt{normal}: the battery does not contain a detected defect;
    \item \texttt{scratched}: visible scratches are present on the casing;
    \item \texttt{swollen}: the battery shows signs of swelling.
\end{itemize}

The model weights have been validated with the four required classes. The live connection between the IMITECH image stream and the Windows inference service remains an integration step. The WPF--Jetson data contract, however, has already been prepared so that the AI service can send a class name, confidence value, inspection identifier, and timestamp without requiring further changes to the robot control interface.

The conveyor transports batteries from the Windows inspection position toward the robot working area. Its speed and the physical distance between the trigger camera and the picking point are used to estimate when the robot should begin its motion. The current interface stores these values in millimetres per second, millimetres, and milliseconds. At the present stage, these parameters describe the conveyor but do not directly control its motor because the conveyor driver and hardware communication protocol are still pending.

The Jetson Nano acts as the control computer for the DOFBOT station. It hosts the Robot API and Vision Trigger modules inside the existing DOFBOT Docker container. A systemd service starts the API automatically after the Jetson boots. The service exposes HTTP endpoints on port 7000 and can be accessed from the Windows workstation through Ethernet or Wi-Fi.

The camera mounted on the DOFBOT has a different role from the IMITECH camera. It does not classify the battery defect. Instead, it monitors a centred Entry Zone and determines when a physical object reaches the robot trigger area. Background subtraction and contour analysis are used to determine whether an object is present. The camera may also estimate the object orientation so that Servo~5 can compensate for changes in battery direction. This separation prevents the robot trigger from depending completely on the AI classification camera.

The DOFBOT contains six controllable servo channels: base, shoulder, elbow, wrist pitch, wrist rotation, and gripper. Robot movements are executed through predefined poses stored in \texttt{robot\_api.py}. Three defect classes are mapped to three different sorting positions, while normal batteries continue to a collection tray at the end of the conveyor without being picked.

The logical system layout can therefore be summarized as follows:

\begin{center}
\texttt{IMITECH camera} $\rightarrow$ \texttt{Windows YOLOv8} $\rightarrow$ \texttt{FIFO classification queue} $\rightarrow$ \texttt{conveyor} $\rightarrow$ \texttt{DOFBOT camera trigger} $\rightarrow$ \texttt{Robot API} $\rightarrow$ \texttt{DOFBOT motion}
\end{center}

For a normal battery, the last stage is replaced by a PASS operation:

\begin{center}
\texttt{Normal result} $\rightarrow$ \texttt{FIFO queue} $\rightarrow$ \texttt{DOFBOT trigger} $\rightarrow$ \texttt{PASS} $\rightarrow$ \texttt{end-of-conveyor tray}
\end{center}

This architecture ensures that the Windows camera determines what the object is, whereas the DOFBOT camera determines when the corresponding action should occur.


\subsubsection{User Interface}
\label{sec:user-interface}

The system is operated through a WPF-based industrial human--machine interface. The interface uses a dark theme, status cards, clearly separated control areas, and responsive scaling. The scaling function allows the same application to be used on the development computer and on a mini workstation with a different display resolution.

The application is divided into three primary tabs: \textbf{AI Camera}, \textbf{DOFBOT Robot \& Camera Control}, and \textbf{Workflow / Conveyor}.

\paragraph{AI Camera tab.}
The AI Camera tab is responsible for the IMITECH camera connected to the Windows computer. The device list is obtained through the MVS SDK. The operator can refresh the list, select the detected camera, start or stop image acquisition, and capture an image.

The live camera feed is displayed in the centre of the page. Information including image resolution, frame rate, and connection status allows the operator to confirm that the stream is operating correctly. The interface also reminds the operator to stop image acquisition in the original MVS application before opening the camera through WPF because only one application should own the stream.

The right side of the tab contains system information and the System Log. The log records important application events, including camera discovery, camera connection, Vision API connection, Robot API status, classification routing, communication errors, and workflow changes. Repeated polling responses are filtered to prevent the same status from generating a large number of duplicate log entries.

The AI readiness indicator states that the trained YOLOv8 weights have been validated, while live inference service integration is still pending. This distinction avoids presenting placeholder data as an actual AI result.

\paragraph{DOFBOT Robot \& Camera Control tab.}
As shown in Figure~\ref{fig:dofbot-robot-camera-control}, the DOFBOT camera and manual robot controls are integrated into a single tab. This layout allows the operator to observe the trigger camera while adjusting robot poses.

The camera section provides fields for the Jetson Robot API address and buttons for connecting, disconnecting, starting monitoring, stopping monitoring, resetting the trigger latch, and moving the robot to the \texttt{VISION\_HOME} pose.

Its status panel displays:

\begin{itemize}
    \item Vision API connection;
    \item vision processing state;
    \item camera-open state;
    \item robot readiness;
    \item automatic trigger state;
    \item trigger-latched state;
    \item processed frame count;
    \item stable detection frame count;
    \item last trigger and job information.
\end{itemize}

The camera viewer displays the Entry Zone used for object-presence detection. When an object is detected, the overlay shows its position, estimated direction, calculated Servo~5 angle, AI class associated with the object, and calculated trigger delay.

The right-side configuration panel contains conveyor timing values, including belt speed, camera-to-pick distance, estimated robot time to reach the gripping point, and processing margin. These parameters allow the system to calculate:

\begin{equation}
    t_{\mathrm{arrival}} = \frac{d_{\mathrm{camera\text{-}to\text{-}pick}}}{v_{\mathrm{belt}}}.
    \label{eq:arrival-time}
\end{equation}

The robot start delay is then obtained by subtracting the robot motion time and processing margin from the estimated arrival time.

The interface also contains an AI classification queue test panel. Until the Windows AI service is connected, the operator can manually enqueue one of the four classes with a confidence value. This test uses the same HTTP endpoint that will later be called by the real AI service. It therefore verifies the integration contract without fabricating a detection inside the final AI result window.

The manual robot control section displays all six servos. Each servo can be adjusted using either a slider or a numerical input field. The displayed angle is updated from the Robot API, allowing movements made through another application to be reflected in the WPF interface when the API can read the current servo positions. Additional buttons provide HOME, VISION HOME, gripper open, gripper close, reset, copy-actual-position, and apply-all functions.

\paragraph{Workflow / Conveyor tab.}
The Workflow / Conveyor tab provides a static process diagram and a live event timeline. The diagram uses swimlanes to show the responsibilities of the Windows inspection station, conveyor, Jetson vision module, and DOFBOT robot.

The four routing branches are explicitly displayed:

\begin{itemize}
    \item \texttt{normal} $\rightarrow$ \texttt{PASS};
    \item \texttt{dented} $\rightarrow$ \texttt{DENTED bin};
    \item \texttt{scratched} $\rightarrow$ \texttt{SCRATCHED bin};
    \item \texttt{swollen} $\rightarrow$ \texttt{SWOLLEN bin}.
\end{itemize}

The Process Information card shows the current workflow stage, latest event, next queued AI result, AI queue size, reference belt speed, automatic-trigger state, and current robot step. This tab provides a high-level operational view without duplicating all detailed camera and servo controls.


\subsubsection{AI Application Processing Flow}
\label{sec:ai-processing-flow}

The AI processing flow begins when the IMITECH camera captures a raw image of a battery. The input image must not contain a manually drawn bounding box, class name, or confidence value. These graphics can introduce artificial visual features and negatively influence model evaluation.

For ordinary prediction, only the raw image is required. For quantitative validation, each raw test image must be accompanied by a corresponding YOLO-format label file containing the ground-truth class and normalized bounding-box coordinates.

The image is passed to the trained YOLOv8 detection model. The current model contains four output classes: \texttt{dented}, \texttt{normal}, \texttt{scratched}, and \texttt{swollen}. During the previous evaluation, the test dataset contained 106 images. The model correctly classified 104 images, while two normal batteries were predicted as swollen. The overall reported values were approximately 0.981 precision, 0.986 recall, 0.981 mAP50, and 0.720 mAP50--95.

The Windows inference service is expected to generate a normalized result containing:

\begin{itemize}
    \item a unique inspection identifier;
    \item predicted class;
    \item confidence value;
    \item bounding-box coordinates;
    \item inference time;
    \item processing timestamp.
\end{itemize}

A default confidence threshold of 0.40 is used by the integration contract. Results below the threshold are rejected and must not produce a robot movement. The maximum number of detections can be limited to one because the intended operating condition contains one battery in each inspection cycle.

After inference, the Windows service sends the result to:

\begin{center}
    \texttt{POST /vision/classifications}
\end{center}

The Vision Trigger module stores accepted results in a first-in, first-out queue. FIFO ordering is important because the IMITECH camera and DOFBOT camera observe the same batteries at different physical positions. The first battery classified by the IMITECH camera should therefore be paired with the first object subsequently detected by the DOFBOT camera.

Each queued result has an expiration time. The default time-to-live is 30 seconds. Expired results are removed to prevent a classification belonging to an old battery from being used for a new object.

When an object remains inside the DOFBOT Entry Zone for the required number of stable frames, the trigger module examines the first valid AI result. One result is consumed for each physical object, including normal batteries.

The routing logic is then applied:

\begin{itemize}
    \item For \texttt{normal}, the API creates a PASS routing event. No robot job is generated.
    \item For \texttt{dented}, a robot job targeting the DENTED location is requested.
    \item For \texttt{scratched}, a robot job targeting the SCRATCHED location is requested.
    \item For \texttt{swollen}, a robot job targeting the SWOLLEN location is requested.
\end{itemize}

For defect classes, the trigger also supplies the estimated Servo~5 angle and start delay. The Robot API places the request into its own job queue. A single robot worker executes the jobs sequentially, preventing two commands from controlling the arm simultaneously.

The defect drop poses are currently marked as uncalibrated. Therefore, the Robot API intentionally rejects automatic defect cycles until the three physical sorting positions have been measured and collision-tested. This is a safety mechanism rather than an incomplete routing implementation.


\subsubsection{Operating Procedures}
\label{sec:operating-procedures}

\paragraph{System start-up.}
Before starting the application, the operator must inspect the robot, conveyor, battery trays, cables, camera mounts, and working area. No person or unrelated object should be inside the DOFBOT workspace.

The Jetson Nano is powered on first. The systemd service should automatically start the Docker container and Robot API. The operator verifies the service using the health endpoint:

\begin{center}
    \texttt{http://<JETSON-IP>:7000/health}
\end{center}

The response should report that the API is operational, the robot is initialized, and the worker thread is alive.

The WPF application is then started on the Windows computer. The operator connects the IMITECH camera from the AI Camera tab. If the camera is currently acquiring images in MVS, acquisition must be stopped before WPF opens the device.

The operator enters the Jetson address in the DOFBOT tab and connects both the Vision API and Robot API. The displayed connection, robot initialization, motion-enabled, and camera-open states must be checked before any movement is commanded.

\paragraph{Monitor-only operation.}
During configuration and camera testing, Auto Trigger must remain disabled. The robot is moved to \texttt{VISION\_HOME}, and DOFBOT camera monitoring is started.

A battery is passed through the Entry Zone to verify that detection becomes true, the trigger latch activates, and the displayed object information changes. Removing the object should allow the latch to rearm after the configured number of clear frames.

The WPF classification test panel can be used to enqueue a class and confidence value. In monitor-only mode, this result can be previewed without automatically executing a robot cycle.

\paragraph{Automatic sorting operation.}
Automatic operation may only begin after the pick poses, three defect drop poses, Servo~5 limits, conveyor timing, and collision clearances have been confirmed.

The Windows AI service is started and connected to the IMITECH image stream. The operator confirms that classifications are entering the FIFO queue in the same order as batteries on the conveyor.

After the robot has reached \texttt{VISION\_HOME}, the conveyor may be started and Auto Trigger enabled. For each battery, the operator monitors the System Log, queued classification, DOFBOT trigger state, and robot job state.

Normal batteries should continue to the tray at the end of the conveyor. Defective batteries should be picked and placed into their respective DENTED, SCRATCHED, or SWOLLEN locations.

\paragraph{Fault and collision response.}
If the robot collides with another object or moves unexpectedly, the operator must immediately stop the robot using the physical power switch or emergency stop mechanism. A software Reset or HOME command must not be used while the robot is physically obstructed.

After power is removed, the cause of the collision must be identified. Possible causes include an incorrect pose, shifted conveyor, wrong battery position, unsuitable Servo~5 correction, mechanical obstruction, or stale classification result.

The System Log and Jetson service journal should be reviewed before restarting. The classification queue and trigger latch should be cleared so that an old AI result cannot be assigned to the next battery.

The robot should initially be tested at a low movement speed and without a battery. HOME or VISION HOME should only be commanded after the workspace is confirmed to be clear.

\paragraph{Shutdown.}
At the end of operation, automatic triggering is disabled first. The conveyor is stopped, vision monitoring is stopped, and the camera streams are closed. If the working area is clear, the robot may be returned to HOME. The WPF application can then be closed, followed by the Jetson Nano and other hardware when required.

At the present commissioning stage, the system should remain in monitor-only or manual test mode until the conveyor driver, live Windows AI inference service, and three defect drop poses have been fully validated.

%
% These headings use \subsubsection so that, when placed under Section 3.2,
% LaTeX can number them automatically as 3.2.4--3.2.7.
% Add \label{fig:dofbot-robot-camera-control} to the Figure 7 environment.

\subsubsection{System Layout}
\label{sec:system-layout}

The proposed battery inspection and sorting system is organized into three main processing areas: the Windows inspection station, the conveyor section, and the Jetson--DOFBOT robot station. This distributed arrangement separates image classification, object transportation, trigger detection, and robot motion control into independent modules. The modular design simplifies testing and allows each subsystem to be developed or replaced without redesigning the entire system.

The Windows inspection station consists of a Windows computer or mini workstation, an IMITECH IMB-770GC industrial camera, the MVS camera driver, the WPF human--machine interface, and the YOLOv8 inference component. The IMITECH camera is positioned above the conveyor to capture detailed images of each battery. Its primary function is to provide high-quality input images for the AI model. Camera communication is handled through the MVS SDK rather than a conventional webcam interface. Therefore, the camera stream must not be opened simultaneously by both the MVS application and the WPF software.

The YOLOv8 model classifies each inspected battery into one of four categories:

\begin{itemize}
    \item \texttt{dented}: the battery casing is deformed or dented;
    \item \texttt{normal}: the battery does not contain a detected defect;
    \item \texttt{scratched}: visible scratches are present on the casing;
    \item \texttt{swollen}: the battery shows signs of swelling.
\end{itemize}

The model weights have been validated with the four required classes. The live connection between the IMITECH image stream and the Windows inference service remains an integration step. The WPF--Jetson data contract, however, has already been prepared so that the AI service can send a class name, confidence value, inspection identifier, and timestamp without requiring further changes to the robot control interface.

The conveyor transports batteries from the Windows inspection position toward the robot working area. Its speed and the physical distance between the trigger camera and the picking point are used to estimate when the robot should begin its motion. The current interface stores these values in millimetres per second, millimetres, and milliseconds. At the present stage, these parameters describe the conveyor but do not directly control its motor because the conveyor driver and hardware communication protocol are still pending.

The Jetson Nano acts as the control computer for the DOFBOT station. It hosts the Robot API and Vision Trigger modules inside the existing DOFBOT Docker container. A systemd service starts the API automatically after the Jetson boots. The service exposes HTTP endpoints on port 7000 and can be accessed from the Windows workstation through Ethernet or Wi-Fi.

The camera mounted on the DOFBOT has a different role from the IMITECH camera. It does not classify the battery defect. Instead, it monitors a centred Entry Zone and determines when a physical object reaches the robot trigger area. Background subtraction and contour analysis are used to determine whether an object is present. The camera may also estimate the object orientation so that Servo~5 can compensate for changes in battery direction. This separation prevents the robot trigger from depending completely on the AI classification camera.

The DOFBOT contains six controllable servo channels: base, shoulder, elbow, wrist pitch, wrist rotation, and gripper. Robot movements are executed through predefined poses stored in \texttt{robot\_api.py}. Three defect classes are mapped to three different sorting positions, while normal batteries continue to a collection tray at the end of the conveyor without being picked.

The logical system layout can therefore be summarized as follows:

\begin{center}
\texttt{IMITECH camera} $\rightarrow$ \texttt{Windows YOLOv8} $\rightarrow$ \texttt{FIFO classification queue} $\rightarrow$ \texttt{conveyor} $\rightarrow$ \texttt{DOFBOT camera trigger} $\rightarrow$ \texttt{Robot API} $\rightarrow$ \texttt{DOFBOT motion}
\end{center}

For a normal battery, the last stage is replaced by a PASS operation:

\begin{center}
\texttt{Normal result} $\rightarrow$ \texttt{FIFO queue} $\rightarrow$ \texttt{DOFBOT trigger} $\rightarrow$ \texttt{PASS} $\rightarrow$ \texttt{end-of-conveyor tray}
\end{center}

This architecture ensures that the Windows camera determines what the object is, whereas the DOFBOT camera determines when the corresponding action should occur.


\subsubsection{User Interface}
\label{sec:user-interface}

The system is operated through a WPF-based industrial human--machine interface. The interface uses a dark theme, status cards, clearly separated control areas, and responsive scaling. The scaling function allows the same application to be used on the development computer and on a mini workstation with a different display resolution.

The application is divided into three primary tabs: \textbf{AI Camera}, \textbf{DOFBOT Robot \& Camera Control}, and \textbf{Workflow / Conveyor}.

\paragraph{AI Camera tab.}
The AI Camera tab is responsible for the IMITECH camera connected to the Windows computer. The device list is obtained through the MVS SDK. The operator can refresh the list, select the detected camera, start or stop image acquisition, and capture an image.

The live camera feed is displayed in the centre of the page. Information including image resolution, frame rate, and connection status allows the operator to confirm that the stream is operating correctly. The interface also reminds the operator to stop image acquisition in the original MVS application before opening the camera through WPF because only one application should own the stream.

The right side of the tab contains system information and the System Log. The log records important application events, including camera discovery, camera connection, Vision API connection, Robot API status, classification routing, communication errors, and workflow changes. Repeated polling responses are filtered to prevent the same status from generating a large number of duplicate log entries.

The AI readiness indicator states that the trained YOLOv8 weights have been validated, while live inference service integration is still pending. This distinction avoids presenting placeholder data as an actual AI result.

\paragraph{DOFBOT Robot \& Camera Control tab.}
As shown in Figure~\ref{fig:dofbot-robot-camera-control}, the DOFBOT camera and manual robot controls are integrated into a single tab. This layout allows the operator to observe the trigger camera while adjusting robot poses.

The camera section provides fields for the Jetson Robot API address and buttons for connecting, disconnecting, starting monitoring, stopping monitoring, resetting the trigger latch, and moving the robot to the \texttt{VISION\_HOME} pose.

Its status panel displays:

\begin{itemize}
    \item Vision API connection;
    \item vision processing state;
    \item camera-open state;
    \item robot readiness;
    \item automatic trigger state;
    \item trigger-latched state;
    \item processed frame count;
    \item stable detection frame count;
    \item last trigger and job information.
\end{itemize}

The camera viewer displays the Entry Zone used for object-presence detection. When an object is detected, the overlay shows its position, estimated direction, calculated Servo~5 angle, AI class associated with the object, and calculated trigger delay.

The right-side configuration panel contains conveyor timing values, including belt speed, camera-to-pick distance, estimated robot time to reach the gripping point, and processing margin. These parameters allow the system to calculate:

\begin{equation}
    t_{\mathrm{arrival}} = \frac{d_{\mathrm{camera\text{-}to\text{-}pick}}}{v_{\mathrm{belt}}}.
    \label{eq:arrival-time}
\end{equation}

The robot start delay is then obtained by subtracting the robot motion time and processing margin from the estimated arrival time.

The interface also contains an AI classification queue test panel. Until the Windows AI service is connected, the operator can manually enqueue one of the four classes with a confidence value. This test uses the same HTTP endpoint that will later be called by the real AI service. It therefore verifies the integration contract without fabricating a detection inside the final AI result window.

The manual robot control section displays all six servos. Each servo can be adjusted using either a slider or a numerical input field. The displayed angle is updated from the Robot API, allowing movements made through another application to be reflected in the WPF interface when the API can read the current servo positions. Additional buttons provide HOME, VISION HOME, gripper open, gripper close, reset, copy-actual-position, and apply-all functions.

\paragraph{Workflow / Conveyor tab.}
The Workflow / Conveyor tab provides a static process diagram and a live event timeline. The diagram uses swimlanes to show the responsibilities of the Windows inspection station, conveyor, Jetson vision module, and DOFBOT robot.

The four routing branches are explicitly displayed:

\begin{itemize}
    \item \texttt{normal} $\rightarrow$ \texttt{PASS};
    \item \texttt{dented} $\rightarrow$ \texttt{DENTED bin};
    \item \texttt{scratched} $\rightarrow$ \texttt{SCRATCHED bin};
    \item \texttt{swollen} $\rightarrow$ \texttt{SWOLLEN bin}.
\end{itemize}

The Process Information card shows the current workflow stage, latest event, next queued AI result, AI queue size, reference belt speed, automatic-trigger state, and current robot step. This tab provides a high-level operational view without duplicating all detailed camera and servo controls.


\subsubsection{AI Application Processing Flow}
\label{sec:ai-processing-flow}

The AI processing flow begins when the IMITECH camera captures a raw image of a battery. The input image must not contain a manually drawn bounding box, class name, or confidence value. These graphics can introduce artificial visual features and negatively influence model evaluation.

For ordinary prediction, only the raw image is required. For quantitative validation, each raw test image must be accompanied by a corresponding YOLO-format label file containing the ground-truth class and normalized bounding-box coordinates.

The image is passed to the trained YOLOv8 detection model. The current model contains four output classes: \texttt{dented}, \texttt{normal}, \texttt{scratched}, and \texttt{swollen}. During the previous evaluation, the test dataset contained 106 images. The model correctly classified 104 images, while two normal batteries were predicted as swollen. The overall reported values were approximately 0.981 precision, 0.986 recall, 0.981 mAP50, and 0.720 mAP50--95.

The Windows inference service is expected to generate a normalized result containing:

\begin{itemize}
    \item a unique inspection identifier;
    \item predicted class;
    \item confidence value;
    \item bounding-box coordinates;
    \item inference time;
    \item processing timestamp.
\end{itemize}

A default confidence threshold of 0.40 is used by the integration contract. Results below the threshold are rejected and must not produce a robot movement. The maximum number of detections can be limited to one because the intended operating condition contains one battery in each inspection cycle.

After inference, the Windows service sends the result to:

\begin{center}
    \texttt{POST /vision/classifications}
\end{center}

The Vision Trigger module stores accepted results in a first-in, first-out queue. FIFO ordering is important because the IMITECH camera and DOFBOT camera observe the same batteries at different physical positions. The first battery classified by the IMITECH camera should therefore be paired with the first object subsequently detected by the DOFBOT camera.

Each queued result has an expiration time. The default time-to-live is 30 seconds. Expired results are removed to prevent a classification belonging to an old battery from being used for a new object.

When an object remains inside the DOFBOT Entry Zone for the required number of stable frames, the trigger module examines the first valid AI result. One result is consumed for each physical object, including normal batteries.

The routing logic is then applied:

\begin{itemize}
    \item For \texttt{normal}, the API creates a PASS routing event. No robot job is generated.
    \item For \texttt{dented}, a robot job targeting the DENTED location is requested.
    \item For \texttt{scratched}, a robot job targeting the SCRATCHED location is requested.
    \item For \texttt{swollen}, a robot job targeting the SWOLLEN location is requested.
\end{itemize}

For defect classes, the trigger also supplies the estimated Servo~5 angle and start delay. The Robot API places the request into its own job queue. A single robot worker executes the jobs sequentially, preventing two commands from controlling the arm simultaneously.

The defect drop poses are currently marked as uncalibrated. Therefore, the Robot API intentionally rejects automatic defect cycles until the three physical sorting positions have been measured and collision-tested. This is a safety mechanism rather than an incomplete routing implementation.


\subsubsection{Operating Procedures}
\label{sec:operating-procedures}

\paragraph{System start-up.}
Before starting the application, the operator must inspect the robot, conveyor, battery trays, cables, camera mounts, and working area. No person or unrelated object should be inside the DOFBOT workspace.

The Jetson Nano is powered on first. The systemd service should automatically start the Docker container and Robot API. The operator verifies the service using the health endpoint:

\begin{center}
    \texttt{http://<JETSON-IP>:7000/health}
\end{center}

The response should report that the API is operational, the robot is initialized, and the worker thread is alive.

The WPF application is then started on the Windows computer. The operator connects the IMITECH camera from the AI Camera tab. If the camera is currently acquiring images in MVS, acquisition must be stopped before WPF opens the device.

The operator enters the Jetson address in the DOFBOT tab and connects both the Vision API and Robot API. The displayed connection, robot initialization, motion-enabled, and camera-open states must be checked before any movement is commanded.

\paragraph{Monitor-only operation.}
During configuration and camera testing, Auto Trigger must remain disabled. The robot is moved to \texttt{VISION\_HOME}, and DOFBOT camera monitoring is started.

A battery is passed through the Entry Zone to verify that detection becomes true, the trigger latch activates, and the displayed object information changes. Removing the object should allow the latch to rearm after the configured number of clear frames.

The WPF classification test panel can be used to enqueue a class and confidence value. In monitor-only mode, this result can be previewed without automatically executing a robot cycle.

\paragraph{Automatic sorting operation.}
Automatic operation may only begin after the pick poses, three defect drop poses, Servo~5 limits, conveyor timing, and collision clearances have been confirmed.

The Windows AI service is started and connected to the IMITECH image stream. The operator confirms that classifications are entering the FIFO queue in the same order as batteries on the conveyor.

After the robot has reached \texttt{VISION\_HOME}, the conveyor may be started and Auto Trigger enabled. For each battery, the operator monitors the System Log, queued classification, DOFBOT trigger state, and robot job state.

Normal batteries should continue to the tray at the end of the conveyor. Defective batteries should be picked and placed into their respective DENTED, SCRATCHED, or SWOLLEN locations.

\paragraph{Fault and collision response.}
If the robot collides with another object or moves unexpectedly, the operator must immediately stop the robot using the physical power switch or emergency stop mechanism. A software Reset or HOME command must not be used while the robot is physically obstructed.

After power is removed, the cause of the collision must be identified. Possible causes include an incorrect pose, shifted conveyor, wrong battery position, unsuitable Servo~5 correction, mechanical obstruction, or stale classification result.

The System Log and Jetson service journal should be reviewed before restarting. The classification queue and trigger latch should be cleared so that an old AI result cannot be assigned to the next battery.

The robot should initially be tested at a low movement speed and without a battery. HOME or VISION HOME should only be commanded after the workspace is confirmed to be clear.

\paragraph{Shutdown.}
At the end of operation, automatic triggering is disabled first. The conveyor is stopped, vision monitoring is stopped, and the camera streams are closed. If the working area is clear, the robot may be returned to HOME. The WPF application can then be closed, followed by the Jetson Nano and other hardware when required.

At the present commissioning stage, the system should remain in monitor-only or manual test mode until the conveyor driver, live Windows AI inference service, and three defect drop poses have been fully validated.

%
% These headings use \subsubsection so that, when placed under Section 3.2,
% LaTeX can number them automatically as 3.2.4--3.2.7.
% Add \label{fig:dofbot-robot-camera-control} to the Figure 7 environment.

\subsubsection{System Layout}
\label{sec:system-layout}

The proposed battery inspection and sorting system is organized into three main processing areas: the Windows inspection station, the conveyor section, and the Jetson--DOFBOT robot station. This distributed arrangement separates image classification, object transportation, trigger detection, and robot motion control into independent modules. The modular design simplifies testing and allows each subsystem to be developed or replaced without redesigning the entire system.

The Windows inspection station consists of a Windows computer or mini workstation, an IMITECH IMB-770GC industrial camera, the MVS camera driver, the WPF human--machine interface, and the YOLOv8 inference component. The IMITECH camera is positioned above the conveyor to capture detailed images of each battery. Its primary function is to provide high-quality input images for the AI model. Camera communication is handled through the MVS SDK rather than a conventional webcam interface. Therefore, the camera stream must not be opened simultaneously by both the MVS application and the WPF software.

The YOLOv8 model classifies each inspected battery into one of four categories:

\begin{itemize}
    \item \texttt{dented}: the battery casing is deformed or dented;
    \item \texttt{normal}: the battery does not contain a detected defect;
    \item \texttt{scratched}: visible scratches are present on the casing;
    \item \texttt{swollen}: the battery shows signs of swelling.
\end{itemize}

The model weights have been validated with the four required classes. The live connection between the IMITECH image stream and the Windows inference service remains an integration step. The WPF--Jetson data contract, however, has already been prepared so that the AI service can send a class name, confidence value, inspection identifier, and timestamp without requiring further changes to the robot control interface.

The conveyor transports batteries from the Windows inspection position toward the robot working area. Its speed and the physical distance between the trigger camera and the picking point are used to estimate when the robot should begin its motion. The current interface stores these values in millimetres per second, millimetres, and milliseconds. At the present stage, these parameters describe the conveyor but do not directly control its motor because the conveyor driver and hardware communication protocol are still pending.

The Jetson Nano acts as the control computer for the DOFBOT station. It hosts the Robot API and Vision Trigger modules inside the existing DOFBOT Docker container. A systemd service starts the API automatically after the Jetson boots. The service exposes HTTP endpoints on port 7000 and can be accessed from the Windows workstation through Ethernet or Wi-Fi.

The camera mounted on the DOFBOT has a different role from the IMITECH camera. It does not classify the battery defect. Instead, it monitors a centred Entry Zone and determines when a physical object reaches the robot trigger area. Background subtraction and contour analysis are used to determine whether an object is present. The camera may also estimate the object orientation so that Servo~5 can compensate for changes in battery direction. This separation prevents the robot trigger from depending completely on the AI classification camera.

The DOFBOT contains six controllable servo channels: base, shoulder, elbow, wrist pitch, wrist rotation, and gripper. Robot movements are executed through predefined poses stored in \texttt{robot\_api.py}. Three defect classes are mapped to three different sorting positions, while normal batteries continue to a collection tray at the end of the conveyor without being picked.

The logical system layout can therefore be summarized as follows:

\begin{center}
\texttt{IMITECH camera} $\rightarrow$ \texttt{Windows YOLOv8} $\rightarrow$ \texttt{FIFO classification queue} $\rightarrow$ \texttt{conveyor} $\rightarrow$ \texttt{DOFBOT camera trigger} $\rightarrow$ \texttt{Robot API} $\rightarrow$ \texttt{DOFBOT motion}
\end{center}

For a normal battery, the last stage is replaced by a PASS operation:

\begin{center}
\texttt{Normal result} $\rightarrow$ \texttt{FIFO queue} $\rightarrow$ \texttt{DOFBOT trigger} $\rightarrow$ \texttt{PASS} $\rightarrow$ \texttt{end-of-conveyor tray}
\end{center}

This architecture ensures that the Windows camera determines what the object is, whereas the DOFBOT camera determines when the corresponding action should occur.


\subsubsection{User Interface}
\label{sec:user-interface}

The system is operated through a WPF-based industrial human--machine interface. The interface uses a dark theme, status cards, clearly separated control areas, and responsive scaling. The scaling function allows the same application to be used on the development computer and on a mini workstation with a different display resolution.

The application is divided into three primary tabs: \textbf{AI Camera}, \textbf{DOFBOT Robot \& Camera Control}, and \textbf{Workflow / Conveyor}.

\paragraph{AI Camera tab.}
The AI Camera tab is responsible for the IMITECH camera connected to the Windows computer. The device list is obtained through the MVS SDK. The operator can refresh the list, select the detected camera, start or stop image acquisition, and capture an image.

The live camera feed is displayed in the centre of the page. Information including image resolution, frame rate, and connection status allows the operator to confirm that the stream is operating correctly. The interface also reminds the operator to stop image acquisition in the original MVS application before opening the camera through WPF because only one application should own the stream.

The right side of the tab contains system information and the System Log. The log records important application events, including camera discovery, camera connection, Vision API connection, Robot API status, classification routing, communication errors, and workflow changes. Repeated polling responses are filtered to prevent the same status from generating a large number of duplicate log entries.

The AI readiness indicator states that the trained YOLOv8 weights have been validated, while live inference service integration is still pending. This distinction avoids presenting placeholder data as an actual AI result.

\paragraph{DOFBOT Robot \& Camera Control tab.}
As shown in Figure~\ref{fig:dofbot-robot-camera-control}, the DOFBOT camera and manual robot controls are integrated into a single tab. This layout allows the operator to observe the trigger camera while adjusting robot poses.

The camera section provides fields for the Jetson Robot API address and buttons for connecting, disconnecting, starting monitoring, stopping monitoring, resetting the trigger latch, and moving the robot to the \texttt{VISION\_HOME} pose.

Its status panel displays:

\begin{itemize}
    \item Vision API connection;
    \item vision processing state;
    \item camera-open state;
    \item robot readiness;
    \item automatic trigger state;
    \item trigger-latched state;
    \item processed frame count;
    \item stable detection frame count;
    \item last trigger and job information.
\end{itemize}

The camera viewer displays the Entry Zone used for object-presence detection. When an object is detected, the overlay shows its position, estimated direction, calculated Servo~5 angle, AI class associated with the object, and calculated trigger delay.

The right-side configuration panel contains conveyor timing values, including belt speed, camera-to-pick distance, estimated robot time to reach the gripping point, and processing margin. These parameters allow the system to calculate:

\begin{equation}
    t_{\mathrm{arrival}} = \frac{d_{\mathrm{camera\text{-}to\text{-}pick}}}{v_{\mathrm{belt}}}.
    \label{eq:arrival-time}
\end{equation}

The robot start delay is then obtained by subtracting the robot motion time and processing margin from the estimated arrival time.

The interface also contains an AI classification queue test panel. Until the Windows AI service is connected, the operator can manually enqueue one of the four classes with a confidence value. This test uses the same HTTP endpoint that will later be called by the real AI service. It therefore verifies the integration contract without fabricating a detection inside the final AI result window.

The manual robot control section displays all six servos. Each servo can be adjusted using either a slider or a numerical input field. The displayed angle is updated from the Robot API, allowing movements made through another application to be reflected in the WPF interface when the API can read the current servo positions. Additional buttons provide HOME, VISION HOME, gripper open, gripper close, reset, copy-actual-position, and apply-all functions.

\paragraph{Workflow / Conveyor tab.}
The Workflow / Conveyor tab provides a static process diagram and a live event timeline. The diagram uses swimlanes to show the responsibilities of the Windows inspection station, conveyor, Jetson vision module, and DOFBOT robot.

The four routing branches are explicitly displayed:

\begin{itemize}
    \item \texttt{normal} $\rightarrow$ \texttt{PASS};
    \item \texttt{dented} $\rightarrow$ \texttt{DENTED bin};
    \item \texttt{scratched} $\rightarrow$ \texttt{SCRATCHED bin};
    \item \texttt{swollen} $\rightarrow$ \texttt{SWOLLEN bin}.
\end{itemize}

The Process Information card shows the current workflow stage, latest event, next queued AI result, AI queue size, reference belt speed, automatic-trigger state, and current robot step. This tab provides a high-level operational view without duplicating all detailed camera and servo controls.


\subsubsection{AI Application Processing Flow}
\label{sec:ai-processing-flow}

The AI processing flow begins when the IMITECH camera captures a raw image of a battery. The input image must not contain a manually drawn bounding box, class name, or confidence value. These graphics can introduce artificial visual features and negatively influence model evaluation.

For ordinary prediction, only the raw image is required. For quantitative validation, each raw test image must be accompanied by a corresponding YOLO-format label file containing the ground-truth class and normalized bounding-box coordinates.

The image is passed to the trained YOLOv8 detection model. The current model contains four output classes: \texttt{dented}, \texttt{normal}, \texttt{scratched}, and \texttt{swollen}. During the previous evaluation, the test dataset contained 106 images. The model correctly classified 104 images, while two normal batteries were predicted as swollen. The overall reported values were approximately 0.981 precision, 0.986 recall, 0.981 mAP50, and 0.720 mAP50--95.

The Windows inference service is expected to generate a normalized result containing:

\begin{itemize}
    \item a unique inspection identifier;
    \item predicted class;
    \item confidence value;
    \item bounding-box coordinates;
    \item inference time;
    \item processing timestamp.
\end{itemize}

A default confidence threshold of 0.40 is used by the integration contract. Results below the threshold are rejected and must not produce a robot movement. The maximum number of detections can be limited to one because the intended operating condition contains one battery in each inspection cycle.

After inference, the Windows service sends the result to:

\begin{center}
    \texttt{POST /vision/classifications}
\end{center}

The Vision Trigger module stores accepted results in a first-in, first-out queue. FIFO ordering is important because the IMITECH camera and DOFBOT camera observe the same batteries at different physical positions. The first battery classified by the IMITECH camera should therefore be paired with the first object subsequently detected by the DOFBOT camera.

Each queued result has an expiration time. The default time-to-live is 30 seconds. Expired results are removed to prevent a classification belonging to an old battery from being used for a new object.

When an object remains inside the DOFBOT Entry Zone for the required number of stable frames, the trigger module examines the first valid AI result. One result is consumed for each physical object, including normal batteries.

The routing logic is then applied:

\begin{itemize}
    \item For \texttt{normal}, the API creates a PASS routing event. No robot job is generated.
    \item For \texttt{dented}, a robot job targeting the DENTED location is requested.
    \item For \texttt{scratched}, a robot job targeting the SCRATCHED location is requested.
    \item For \texttt{swollen}, a robot job targeting the SWOLLEN location is requested.
\end{itemize}

For defect classes, the trigger also supplies the estimated Servo~5 angle and start delay. The Robot API places the request into its own job queue. A single robot worker executes the jobs sequentially, preventing two commands from controlling the arm simultaneously.

The defect drop poses are currently marked as uncalibrated. Therefore, the Robot API intentionally rejects automatic defect cycles until the three physical sorting positions have been measured and collision-tested. This is a safety mechanism rather than an incomplete routing implementation.


\subsubsection{Operating Procedures}
\label{sec:operating-procedures}

\paragraph{System start-up.}
Before starting the application, the operator must inspect the robot, conveyor, battery trays, cables, camera mounts, and working area. No person or unrelated object should be inside the DOFBOT workspace.

The Jetson Nano is powered on first. The systemd service should automatically start the Docker container and Robot API. The operator verifies the service using the health endpoint:

\begin{center}
    \texttt{http://<JETSON-IP>:7000/health}
\end{center}

The response should report that the API is operational, the robot is initialized, and the worker thread is alive.

The WPF application is then started on the Windows computer. The operator connects the IMITECH camera from the AI Camera tab. If the camera is currently acquiring images in MVS, acquisition must be stopped before WPF opens the device.

The operator enters the Jetson address in the DOFBOT tab and connects both the Vision API and Robot API. The displayed connection, robot initialization, motion-enabled, and camera-open states must be checked before any movement is commanded.

\paragraph{Monitor-only operation.}
During configuration and camera testing, Auto Trigger must remain disabled. The robot is moved to \texttt{VISION\_HOME}, and DOFBOT camera monitoring is started.

A battery is passed through the Entry Zone to verify that detection becomes true, the trigger latch activates, and the displayed object information changes. Removing the object should allow the latch to rearm after the configured number of clear frames.

The WPF classification test panel can be used to enqueue a class and confidence value. In monitor-only mode, this result can be previewed without automatically executing a robot cycle.

\paragraph{Automatic sorting operation.}
Automatic operation may only begin after the pick poses, three defect drop poses, Servo~5 limits, conveyor timing, and collision clearances have been confirmed.

The Windows AI service is started and connected to the IMITECH image stream. The operator confirms that classifications are entering the FIFO queue in the same order as batteries on the conveyor.

After the robot has reached \texttt{VISION\_HOME}, the conveyor may be started and Auto Trigger enabled. For each battery, the operator monitors the System Log, queued classification, DOFBOT trigger state, and robot job state.

Normal batteries should continue to the tray at the end of the conveyor. Defective batteries should be picked and placed into their respective DENTED, SCRATCHED, or SWOLLEN locations.

\paragraph{Fault and collision response.}
If the robot collides with another object or moves unexpectedly, the operator must immediately stop the robot using the physical power switch or emergency stop mechanism. A software Reset or HOME command must not be used while the robot is physically obstructed.

After power is removed, the cause of the collision must be identified. Possible causes include an incorrect pose, shifted conveyor, wrong battery position, unsuitable Servo~5 correction, mechanical obstruction, or stale classification result.

The System Log and Jetson service journal should be reviewed before restarting. The classification queue and trigger latch should be cleared so that an old AI result cannot be assigned to the next battery.

The robot should initially be tested at a low movement speed and without a battery. HOME or VISION HOME should only be commanded after the workspace is confirmed to be clear.

\paragraph{Shutdown.}
At the end of operation, automatic triggering is disabled first. The conveyor is stopped, vision monitoring is stopped, and the camera streams are closed. If the working area is clear, the robot may be returned to HOME. The WPF application can then be closed, followed by the Jetson Nano and other hardware when required.

At the present commissioning stage, the system should remain in monitor-only or manual test mode until the conveyor driver, live Windows AI inference service, and three defect drop poses have been fully validated.

%
% These headings use \subsubsection so that, when placed under Section 3.2,
% LaTeX can number them automatically as 3.2.4--3.2.7.
% Add \label{fig:dofbot-robot-camera-control} to the Figure 7 environment.

\subsubsection{System Layout}
\label{sec:system-layout}

The proposed battery inspection and sorting system is organized into three main processing areas: the Windows inspection station, the conveyor section, and the Jetson--DOFBOT robot station. This distributed arrangement separates image classification, object transportation, trigger detection, and robot motion control into independent modules. The modular design simplifies testing and allows each subsystem to be developed or replaced without redesigning the entire system.

The Windows inspection station consists of a Windows computer or mini workstation, an IMITECH IMB-770GC industrial camera, the MVS camera driver, the WPF human--machine interface, and the YOLOv8 inference component. The IMITECH camera is positioned above the conveyor to capture detailed images of each battery. Its primary function is to provide high-quality input images for the AI model. Camera communication is handled through the MVS SDK rather than a conventional webcam interface. Therefore, the camera stream must not be opened simultaneously by both the MVS application and the WPF software.

The YOLOv8 model classifies each inspected battery into one of four categories:

\begin{itemize}
    \item \texttt{dented}: the battery casing is deformed or dented;
    \item \texttt{normal}: the battery does not contain a detected defect;
    \item \texttt{scratched}: visible scratches are present on the casing;
    \item \texttt{swollen}: the battery shows signs of swelling.
\end{itemize}

The model weights have been validated with the four required classes. The live connection between the IMITECH image stream and the Windows inference service remains an integration step. The WPF--Jetson data contract, however, has already been prepared so that the AI service can send a class name, confidence value, inspection identifier, and timestamp without requiring further changes to the robot control interface.

The conveyor transports batteries from the Windows inspection position toward the robot working area. Its speed and the physical distance between the trigger camera and the picking point are used to estimate when the robot should begin its motion. The current interface stores these values in millimetres per second, millimetres, and milliseconds. At the present stage, these parameters describe the conveyor but do not directly control its motor because the conveyor driver and hardware communication protocol are still pending.

The Jetson Nano acts as the control computer for the DOFBOT station. It hosts the Robot API and Vision Trigger modules inside the existing DOFBOT Docker container. A systemd service starts the API automatically after the Jetson boots. The service exposes HTTP endpoints on port 7000 and can be accessed from the Windows workstation through Ethernet or Wi-Fi.

The camera mounted on the DOFBOT has a different role from the IMITECH camera. It does not classify the battery defect. Instead, it monitors a centred Entry Zone and determines when a physical object reaches the robot trigger area. Background subtraction and contour analysis are used to determine whether an object is present. The camera may also estimate the object orientation so that Servo~5 can compensate for changes in battery direction. This separation prevents the robot trigger from depending completely on the AI classification camera.

The DOFBOT contains six controllable servo channels: base, shoulder, elbow, wrist pitch, wrist rotation, and gripper. Robot movements are executed through predefined poses stored in \texttt{robot\_api.py}. Three defect classes are mapped to three different sorting positions, while normal batteries continue to a collection tray at the end of the conveyor without being picked.

The logical system layout can therefore be summarized as follows:

\begin{center}
\texttt{IMITECH camera} $\rightarrow$ \texttt{Windows YOLOv8} $\rightarrow$ \texttt{FIFO classification queue} $\rightarrow$ \texttt{conveyor} $\rightarrow$ \texttt{DOFBOT camera trigger} $\rightarrow$ \texttt{Robot API} $\rightarrow$ \texttt{DOFBOT motion}
\end{center}

For a normal battery, the last stage is replaced by a PASS operation:

\begin{center}
\texttt{Normal result} $\rightarrow$ \texttt{FIFO queue} $\rightarrow$ \texttt{DOFBOT trigger} $\rightarrow$ \texttt{PASS} $\rightarrow$ \texttt{end-of-conveyor tray}
\end{center}

This architecture ensures that the Windows camera determines what the object is, whereas the DOFBOT camera determines when the corresponding action should occur.


\subsubsection{User Interface}
\label{sec:user-interface}

The system is operated through a WPF-based industrial human--machine interface. The interface uses a dark theme, status cards, clearly separated control areas, and responsive scaling. The scaling function allows the same application to be used on the development computer and on a mini workstation with a different display resolution.

The application is divided into three primary tabs: \textbf{AI Camera}, \textbf{DOFBOT Robot \& Camera Control}, and \textbf{Workflow / Conveyor}.

\paragraph{AI Camera tab.}
The AI Camera tab is responsible for the IMITECH camera connected to the Windows computer. The device list is obtained through the MVS SDK. The operator can refresh the list, select the detected camera, start or stop image acquisition, and capture an image.

The live camera feed is displayed in the centre of the page. Information including image resolution, frame rate, and connection status allows the operator to confirm that the stream is operating correctly. The interface also reminds the operator to stop image acquisition in the original MVS application before opening the camera through WPF because only one application should own the stream.

The right side of the tab contains system information and the System Log. The log records important application events, including camera discovery, camera connection, Vision API connection, Robot API status, classification routing, communication errors, and workflow changes. Repeated polling responses are filtered to prevent the same status from generating a large number of duplicate log entries.

The AI readiness indicator states that the trained YOLOv8 weights have been validated, while live inference service integration is still pending. This distinction avoids presenting placeholder data as an actual AI result.

\paragraph{DOFBOT Robot \& Camera Control tab.}
As shown in Figure~\ref{fig:dofbot-robot-camera-control}, the DOFBOT camera and manual robot controls are integrated into a single tab. This layout allows the operator to observe the trigger camera while adjusting robot poses.

The camera section provides fields for the Jetson Robot API address and buttons for connecting, disconnecting, starting monitoring, stopping monitoring, resetting the trigger latch, and moving the robot to the \texttt{VISION\_HOME} pose.

Its status panel displays:

\begin{itemize}
    \item Vision API connection;
    \item vision processing state;
    \item camera-open state;
    \item robot readiness;
    \item automatic trigger state;
    \item trigger-latched state;
    \item processed frame count;
    \item stable detection frame count;
    \item last trigger and job information.
\end{itemize}

The camera viewer displays the Entry Zone used for object-presence detection. When an object is detected, the overlay shows its position, estimated direction, calculated Servo~5 angle, AI class associated with the object, and calculated trigger delay.

The right-side configuration panel contains conveyor timing values, including belt speed, camera-to-pick distance, estimated robot time to reach the gripping point, and processing margin. These parameters allow the system to calculate:

\begin{equation}
    t_{\mathrm{arrival}} = \frac{d_{\mathrm{camera\text{-}to\text{-}pick}}}{v_{\mathrm{belt}}}.
    \label{eq:arrival-time}
\end{equation}

The robot start delay is then obtained by subtracting the robot motion time and processing margin from the estimated arrival time.

The interface also contains an AI classification queue test panel. Until the Windows AI service is connected, the operator can manually enqueue one of the four classes with a confidence value. This test uses the same HTTP endpoint that will later be called by the real AI service. It therefore verifies the integration contract without fabricating a detection inside the final AI result window.

The manual robot control section displays all six servos. Each servo can be adjusted using either a slider or a numerical input field. The displayed angle is updated from the Robot API, allowing movements made through another application to be reflected in the WPF interface when the API can read the current servo positions. Additional buttons provide HOME, VISION HOME, gripper open, gripper close, reset, copy-actual-position, and apply-all functions.

\paragraph{Workflow / Conveyor tab.}
The Workflow / Conveyor tab provides a static process diagram and a live event timeline. The diagram uses swimlanes to show the responsibilities of the Windows inspection station, conveyor, Jetson vision module, and DOFBOT robot.

The four routing branches are explicitly displayed:

\begin{itemize}
    \item \texttt{normal} $\rightarrow$ \texttt{PASS};
    \item \texttt{dented} $\rightarrow$ \texttt{DENTED bin};
    \item \texttt{scratched} $\rightarrow$ \texttt{SCRATCHED bin};
    \item \texttt{swollen} $\rightarrow$ \texttt{SWOLLEN bin}.
\end{itemize}

The Process Information card shows the current workflow stage, latest event, next queued AI result, AI queue size, reference belt speed, automatic-trigger state, and current robot step. This tab provides a high-level operational view without duplicating all detailed camera and servo controls.


\subsubsection{AI Application Processing Flow}
\label{sec:ai-processing-flow}

The AI processing flow begins when the IMITECH camera captures a raw image of a battery. The input image must not contain a manually drawn bounding box, class name, or confidence value. These graphics can introduce artificial visual features and negatively influence model evaluation.

For ordinary prediction, only the raw image is required. For quantitative validation, each raw test image must be accompanied by a corresponding YOLO-format label file containing the ground-truth class and normalized bounding-box coordinates.

The image is passed to the trained YOLOv8 detection model. The current model contains four output classes: \texttt{dented}, \texttt{normal}, \texttt{scratched}, and \texttt{swollen}. During the previous evaluation, the test dataset contained 106 images. The model correctly classified 104 images, while two normal batteries were predicted as swollen. The overall reported values were approximately 0.981 precision, 0.986 recall, 0.981 mAP50, and 0.720 mAP50--95.

The Windows inference service is expected to generate a normalized result containing:

\begin{itemize}
    \item a unique inspection identifier;
    \item predicted class;
    \item confidence value;
    \item bounding-box coordinates;
    \item inference time;
    \item processing timestamp.
\end{itemize}

A default confidence threshold of 0.40 is used by the integration contract. Results below the threshold are rejected and must not produce a robot movement. The maximum number of detections can be limited to one because the intended operating condition contains one battery in each inspection cycle.

After inference, the Windows service sends the result to:

\begin{center}
    \texttt{POST /vision/classifications}
\end{center}

The Vision Trigger module stores accepted results in a first-in, first-out queue. FIFO ordering is important because the IMITECH camera and DOFBOT camera observe the same batteries at different physical positions. The first battery classified by the IMITECH camera should therefore be paired with the first object subsequently detected by the DOFBOT camera.

Each queued result has an expiration time. The default time-to-live is 30 seconds. Expired results are removed to prevent a classification belonging to an old battery from being used for a new object.

When an object remains inside the DOFBOT Entry Zone for the required number of stable frames, the trigger module examines the first valid AI result. One result is consumed for each physical object, including normal batteries.

The routing logic is then applied:

\begin{itemize}
    \item For \texttt{normal}, the API creates a PASS routing event. No robot job is generated.
    \item For \texttt{dented}, a robot job targeting the DENTED location is requested.
    \item For \texttt{scratched}, a robot job targeting the SCRATCHED location is requested.
    \item For \texttt{swollen}, a robot job targeting the SWOLLEN location is requested.
\end{itemize}

For defect classes, the trigger also supplies the estimated Servo~5 angle and start delay. The Robot API places the request into its own job queue. A single robot worker executes the jobs sequentially, preventing two commands from controlling the arm simultaneously.

The defect drop poses are currently marked as uncalibrated. Therefore, the Robot API intentionally rejects automatic defect cycles until the three physical sorting positions have been measured and collision-tested. This is a safety mechanism rather than an incomplete routing implementation.


\subsubsection{Operating Procedures}
\label{sec:operating-procedures}

\paragraph{System start-up.}
Before starting the application, the operator must inspect the robot, conveyor, battery trays, cables, camera mounts, and working area. No person or unrelated object should be inside the DOFBOT workspace.

The Jetson Nano is powered on first. The systemd service should automatically start the Docker container and Robot API. The operator verifies the service using the health endpoint:

\begin{center}
    \texttt{http://<JETSON-IP>:7000/health}
\end{center}

The response should report that the API is operational, the robot is initialized, and the worker thread is alive.

The WPF application is then started on the Windows computer. The operator connects the IMITECH camera from the AI Camera tab. If the camera is currently acquiring images in MVS, acquisition must be stopped before WPF opens the device.

The operator enters the Jetson address in the DOFBOT tab and connects both the Vision API and Robot API. The displayed connection, robot initialization, motion-enabled, and camera-open states must be checked before any movement is commanded.

\paragraph{Monitor-only operation.}
During configuration and camera testing, Auto Trigger must remain disabled. The robot is moved to \texttt{VISION\_HOME}, and DOFBOT camera monitoring is started.

A battery is passed through the Entry Zone to verify that detection becomes true, the trigger latch activates, and the displayed object information changes. Removing the object should allow the latch to rearm after the configured number of clear frames.

The WPF classification test panel can be used to enqueue a class and confidence value. In monitor-only mode, this result can be previewed without automatically executing a robot cycle.

\paragraph{Automatic sorting operation.}
Automatic operation may only begin after the pick poses, three defect drop poses, Servo~5 limits, conveyor timing, and collision clearances have been confirmed.

The Windows AI service is started and connected to the IMITECH image stream. The operator confirms that classifications are entering the FIFO queue in the same order as batteries on the conveyor.

After the robot has reached \texttt{VISION\_HOME}, the conveyor may be started and Auto Trigger enabled. For each battery, the operator monitors the System Log, queued classification, DOFBOT trigger state, and robot job state.

Normal batteries should continue to the tray at the end of the conveyor. Defective batteries should be picked and placed into their respective DENTED, SCRATCHED, or SWOLLEN locations.

\paragraph{Fault and collision response.}
If the robot collides with another object or moves unexpectedly, the operator must immediately stop the robot using the physical power switch or emergency stop mechanism. A software Reset or HOME command must not be used while the robot is physically obstructed.

After power is removed, the cause of the collision must be identified. Possible causes include an incorrect pose, shifted conveyor, wrong battery position, unsuitable Servo~5 correction, mechanical obstruction, or stale classification result.

The System Log and Jetson service journal should be reviewed before restarting. The classification queue and trigger latch should be cleared so that an old AI result cannot be assigned to the next battery.

The robot should initially be tested at a low movement speed and without a battery. HOME or VISION HOME should only be commanded after the workspace is confirmed to be clear.

\paragraph{Shutdown.}
At the end of operation, automatic triggering is disabled first. The conveyor is stopped, vision monitoring is stopped, and the camera streams are closed. If the working area is clear, the robot may be returned to HOME. The WPF application can then be closed, followed by the Jetson Nano and other hardware when required.

At the present commissioning stage, the system should remain in monitor-only or manual test mode until the conveyor driver, live Windows AI inference service, and three defect drop poses have been fully validated.
